\section{导言：数据不是从天上掉下来的}

上一讲已经把数据来源、抓取和清洗的大图景铺开了。这一讲要回答的不是“数据从哪里来”，而是更具体的三件事：
\textbf{怎么从海量原始数据里挑出像目标数据的子集，怎么把这套机制复用到不同过滤任务里，以及怎么在近重复规模下保持线性级别的效率。}

\begin{importantbox}{本讲的判断标准}
\begin{enumerate}
    \item 结果要足够像目标数据，但不能只是复制目标数据。
    \item 方法要足够快，能在网页级别数据上跑完。
    \item 方法要足够粗糙，粗糙到可以规模化，但又不能粗糙到失去控制。
\end{enumerate}
\end{importantbox}

\begin{figure}[H]
\centering
\resizebox{\textwidth}{!}{%
\begin{tikzpicture}[node distance=1.4cm, every node/.style={font=\small}, scale=0.88, transform shape]
\node[draw, rounded corners, fill=blue!10, minimum width=3cm, minimum height=0.9cm] (live) {live service};
\node[draw, rounded corners, fill=blue!12, right=of live, minimum width=3cm, minimum height=0.9cm] (crawl) {crawl / dump};
\node[draw, rounded corners, fill=blue!16, right=of crawl, minimum width=3.2cm, minimum height=0.9cm] (html) {HTML $\to$ text};
\node[draw, rounded corners, fill=blue!20, right=of html, minimum width=3.3cm, minimum height=0.9cm] (filter) {filter \& dedup};
\node[draw, rounded corners, fill=blue!24, right=of filter, minimum width=3.4cm, minimum height=0.9cm] (train) {training corpus};
\draw[-{Latex[length=3mm]}, thick] (live) -- (crawl);
\draw[-{Latex[length=3mm]}, thick] (crawl) -- (html);
\draw[-{Latex[length=3mm]}, thick] (html) -- (filter);
\draw[-{Latex[length=3mm]}, thick] (filter) -- (train);
\end{tikzpicture}%
}
\caption{从活跃服务到训练语料\protect\footnotemark}
\end{figure}
\footnotetext{字幕 00:00:05--00:01:25。讲者强调数据来自 live services，之后要经历 crawl/dump、HTML 转文本、过滤与去重。}

\begin{figure}[H]
\centering
\begin{tikzpicture}[every node/.style={font=\small, align=center}, scale=0.88, transform shape]
\node[draw, rounded corners, fill=green!10, minimum width=2.8cm, minimum height=0.95cm] (target) at (0,0) {目标数据 $T$\\少而好};
\node[draw, rounded corners, fill=orange!14, minimum width=3cm, minimum height=0.95cm] (raw) at (4.4,0) {原始数据 $R$\\多而杂};
\node[draw, rounded corners, fill=yellow!18, minimum width=3.2cm, minimum height=0.95cm] (score) at (8.9,0) {打分函数 $s(x)$\\perplexity / classifier / weight};
\node[draw, rounded corners, fill=blue!14, minimum width=3.3cm, minimum height=0.95cm] (subset) at (13.5,0) {子集 $R'$\\保留或重采样};
\draw[-{Latex[length=3mm]}, thick] (target) -- (score);
\draw[-{Latex[length=3mm]}, thick] (raw) -- (score);
\draw[-{Latex[length=3mm]}, thick] (score) -- (subset);
\end{tikzpicture}
\caption{过滤问题的抽象：从少量目标数据和大量原始数据里学出一个可打分的选择器}
\end{figure}

\begin{knowledgebox}{为什么这个抽象很通用}
KenLM、fastText、DSIR、语言识别、质量过滤、毒性过滤、近重复召回，本质上都可以塞进同一个框架：先学一个粗糙但快的打分器，再根据分数决定保留哪些样本。
\end{knowledgebox}

